% !TeX root = ../papers/nc-p3.tex
\section{孤立した直線ファイバーを許す全曲面枝}
\label{sec:isolated-surface-fibers}
\begin{lemma}\label{lem:plane-transition-rigidity}
標準正則窓で、各位置の面の一般点を前後に対応させる双有理な整曲面の列があるとする。
一面が平面なら、全てが平面で全写像は射影同型である。別の成分や特殊ファイバーに仮定はない。
\end{lemma}
\begin{proof}
三次元射影は既に可換なので除く。各面は平面か整二次曲面である。
位置0の平面の方程式を $a$ とする。$a$ に始まる純関係がないときは $R_0$ の平面制限が六次元である。
前節の三文字倍と一次関係の証明により、次の面は平面で、対応は一次になる。
そうでない辺があれば、従って $a$ に始まる純関係があり、その直前の $R_{-1}$ に
$a$ に終わる純関係はない。今の議論を逆順に適用すれば、一つ手前は平面で平面0への対応は射影同型である。

前の平面と今の平面の一次式を同じ三空間 $X$ と同定する。
四テンソル $w$ のこの二位置を両平面に制限したものを $w'$ とする。残る二位置を $Y=V_1,Z=V_2$ と書く。
平面グラフ上で消えるので
\[
 w'\in\Lambda^2 X\otimes Y\otimes Z .
\]
しかも $(X,X)\mid(Y,Z)$ の階数は3である。実際、前の辺をまず平面0に射影した段階では
六関係は独立で、次に前の平面へ制限して失う次元は高々 $\dim X=3$ である。
またこの切断の $Y,Z$ 側の像は $R_1$ に入る。

体積形を一つ固定すると、各 $z\in Z^*$ の縮約は線型写像
$T_z:Y^*\to X^*$ を与える。次の面への双有理写像を $f$ とすると、
一般 $x\in\mathbb P(X^*)$ に対して $T_z f(x)$ は $x$ に比例する（零を含む）。
実際 $w'$ は $x,f(x)$ を代入すれば前位置の任意の値に対して消え、三次の交代行列は
その対応するベクトルを核に持つ。像面が平面のとき、ある $T_z$ がそこに非零に制限されれば、
逆写像はこの線型写像であり、平面間の射影同型である。全ての $T_z$ が像平面上で零なら、
$w'$ の $Y$ 支持は平面方程式一本で、切断の階数3により同じ最初の文字の純関係三つを得て矛盾する。

像面が整二次曲面 $Q$ なら、非零の $T_z$ は全て階数3で、その射影化は $Q$ 上で同じ逆写像を与える。
このような線型写像 $T$ は定数倍を除き一つしかない。実際、一つを座標 $y=(u,t)$ で $T y=u$ とし、
他を $T'y=A u+t b$ とする。$u\times(Au+t b)$ の各座標は $Q$ の二次方程式の定数倍である。
それらと $u$ の内積が恒等零であるため、その定数ベクトルは零である。
$t$ の係数から $b=0$、全ての $u$ が $A$ の固有ベクトルなので $A$ はスカラーである。
従って $w'$ の $Z$ 支持が一本となり、今度は同じ末尾文字の三純関係を生じる。
これは定理\ref{thm:branching-centers-curve}に反する。以上を各位置と逆向きにも適用する。
\end{proof}

\begin{lemma}\label{lem:plane-koszul-isolated-exact}
平面の射影同型の四辺続きに沿った標準Koszul複体が、平面の有限個の点を除いて完全なら、
その平面上の層複体は全点で局所分裂完全である。
\end{lemma}
\begin{proof}
各平面の座標を同一に取ると、複体は
\[
 0\longrightarrow\mathcal O(-4)\longrightarrow 4\mathcal O(-3)
 \longrightarrow 6\mathcal O(-2)\longrightarrow 4\mathcal O(-1)\longrightarrow\mathcal O\longrightarrow0
\]
であり、両端の写像は点の非零評価から部分束への単射、全射となる。
最初の余核を $U$、最後の核を $W$ として $U\to6\mathcal O(-2)\to W$ を得る。
後半の核 $K$ は局所自由である。各閉点の二次元正則局所環で後半の像は自由加群の部分加群なので
深さ1以上であり、深さ補題から核は深さ2、従って自由である。
$U\to K$ は同階数で有限点を除いて同型である。その行列式の零点があれば因子を含むから、全点で同型である。
残る唯一のホモロジーは右の有限長余核 $D$ である。一方で捩り後のEuler標数の交代和は
平面の二次Hilbert多項式の四階差分で零なので $\operatorname{length}D=0$ である。
これで完全であり、各核が束であるためファイバーでも完全となる。
\end{proof}

\begin{theorem}\label{thm:isolated-surface-fibers}
標準正則窓の全射影で、直線ファイバーの中心が有限個であるとする。
いずれかの射影像が二次元以上なら、窓は可換多項式、正規一次直線列による平面商、
滑らかな二次曲面商、整二次錐商の既出の添字保持分類に入る。純関係の有無は仮定しない。
\end{theorem}
\begin{proof}
三次元像を除き、系\ref{cor:pure-free-surface-complete}の冒頭と同様に、
面の一般点の相手は再び面の一般点で、一点対の直交不等式から前後に延長できる。
ここで必要なのは両方向の直線ファイバーの中心が有限であることだけである。
平面がなければ定理\ref{thm:birational-quadric-chains}を使う。
平面があれば補題\ref{lem:plane-transition-rigidity}により、全位置の平面商の座標を三変数多項式の反復に取れる。
この商をテンソルした単純加群の線型解消を平面上に層化する。中央二写像は隣接関係の点評価と相補関係の
点評価の転置であり、各平面で有限点以外は両方階数3である。従って補題\ref{lem:plane-koszul-isolated-exact}
により平面上で完全である。定理\ref{thm:plane-surface-normal}の有限長ホモロジーと両側極小解消の議論を
そのまま適用すると、一次核は全位置で正則正規となり、平面商の分類を得る。
\end{proof}
